\section*{Abstract}
In vielen praktischen Anwendungen liegen strukturierte und unstrukturierte Daten gleichzeitig vor, deren effektive gemeinsame Nutzung weiterhin eine zentrale Herausforderung im maschinellen Lernen darstellt. Einfache Fusionsstrategien wie Addition oder Konkatenation erlauben oft nur eingeschränkte Interaktionen zwischen Text- und Tabellendaten. Ziel dieser Arbeit ist es daher zu untersuchen, ob Cross-Attention-Mechanismen eine leistungsfähigere Alternative für die multimodale Feature-Fusion bieten.

Dazu wird die Transformer-basierte Architektur CrossBERT entwickelt, die beide Modalitäten mittels Cross-Attention verknüpft. Drei Varianten (Early, Late, Hybrid) werden auf einem realen Datensatz aus dem Versicherungsbereich sowie zwei öffentlichen Benchmarks implementiert und evaluiert. Im Vergleich zu etablierten Baselines erfolgt die Bewertung anhand gängiger Klassifikationsmetriken und Ablationsstudien.

Die Ergebnisse zeigen auf dem primären Datensatz moderate Zugewinne in der Average Precision bei vergleichbaren AUC-ROC-Werten. Auf einem externen Benchmark mit starken multimodalen Abhängigkeiten erzielt insbesondere die hybride Variante deutliche Verbesserungen, während sich auf einem stark textdominierten Datensatz ein \enquote{Ceiling-Effekt} zeigt. Diese Befunde verdeutlichen, dass der Mehrwert komplexer Fusionsmechanismen maßgeblich von der Datencharakteristik abhängt.

Insgesamt erweist sich Cross-Attention als besonders sinnvoll, wenn beide Modalitäten komplementäre Informationen liefern. Die vorgestellte Architektur bietet eine robuste Grundlage für multimodale Anwendungen, erfordert jedoch eine sorgfältige Abwägung zwischen Leistungsgewinn, Modellkomplexität und Einsatzszenario.
